% Generated draft: requires mathematical and source review.
\documentclass[11pt]{article}
\usepackage[margin=25mm]{geometry}
\usepackage{fontspec}
\setmainfont{Latin Modern Roman}
\usepackage{amsmath,amssymb,mathtools,unicode-math}
\setmathfont{Latin Modern Math}
\usepackage{enumitem,needspace,xurl,hyperref}
\setlength{\parindent}{0pt}
\setlength{\parskip}{5pt}
\setlength{\emergencystretch}{3em}
\newcommand{\N}{\mathbb{N}}
\newcommand{\Z}{\mathbb{Z}}
\newcommand{\Q}{\mathbb{Q}}
\newcommand{\R}{\mathbb{R}}
\newcommand{\C}{\mathbb{C}}
\newcommand{\F}{\mathbb{F}}
\newcommand{\A}{\mathbb{A}}
\newcommand{\PP}{\mathbb{P}}
\newcommand{\E}{\mathbb{E}}
\newcommand{\eps}{\varepsilon}
\newcommand{\dd}{\,\mathrm{d}}
\begin{document}
\section*{UnsolvedMath ID: 30005846}
\textbf{Source problem code:} catalogue entry
\subsection{Definitions and mathematical statement}
For a specified family \(\mathcal S\) of groups and a specified equivalence relation \(\sim\), a complete invariant is an assignment \(I\colon\mathcal S\to\mathcal P\) to a parameter collection \(\mathcal P\) such that
\[
\forall G,H\in\mathcal S,\qquad G\sim H\ \Longleftrightarrow\ I(G)=I(H).
\]
A classification must also specify which parameter values occur.

The supplied problem requests: \emph{Classify Stein's groups as topological full groups.} In the preceding terminology, this asks for a description of the relevant family \(\mathcal S\), its intended equivalence relation, and complete invariants, using the topological full group viewpoint.

The record does not define Stein's groups, their admissible parameters, or the dynamical systems or groupoids whose topological full groups are intended. It also does not specify whether classification concerns abstract group isomorphism, conjugacy of actions, isomorphism of groupoids, or another equivalence relation. Thus the displayed criterion is a classification template, not a fully specified conjecture; neither \(\mathcal S\) nor \(\sim\) can be fixed from the supplied evidence.

The supplied dated assessment reports a classification theorem under the hypotheses that the slope groups are finitely generated and the additive groups have rank at least two, while describing the unrestricted classification as incomplete. These are qualifications of the supplied literature assessment, not additional hypotheses silently imposed on the original problem. The record does not supply the theorem's full parameter conditions or conclusion.
\subsection{Short English statement}
The problem asks for a classification of Stein's groups through their description as topological full groups. The source does not specify the family, its parameters, or the equivalence relation to classify, so a precise unconditional formulation requires further source clarification. The supplied dated curation labels the problem partially solved and reports a classification under additional hypotheses.
\subsection{Sources}
\begin{itemize}
\item Supplied local UnsolvedMath registry record, permanent ID \(30005846\), title \emph{Classification of Stein Groups}, source problem code \texttt{catalogue entry}, in the Oberwolfach Reports Open Problems collection.
\item \emph{Combinatorial \(*\)-algebras} (2024), Oberwolfach Reports, pp.~731--788, DOI \texttt{10.4171/owr/2024/12}. Supplied report link: \url{https://doi.org/10.4171/owr/2024/12}. The record associates the text \emph{Stein's groups as topological full groups} with Owen Tanner; it does not explicitly identify a problem proposer or provide a precise page for the problem.
\item \emph{Classifying Stein's groups}, supplied literature reference: \url{https://arxiv.org/abs/2412.05492}.
\item \emph{Topological full groups thesis}, supplied literature reference: \url{https://theses.gla.ac.uk/84586/2/2024TannerPhD.pdf}.
\item Literature assessment dated exactly \texttt{2026-08-22}: the supplied curation labels the problem \texttt{partially\_solved} and reports a classification for finitely generated slope groups and additive groups of rank at least two, with the unrestricted classification incomplete. This is supplied dated curation, not an independently verified assessment. The background explicitly requests independent verification of time-sensitive status and cited claims.
\item The registry labels the statement \texttt{corrected\_verified} and explains that only formula delimiters, notation, grammar, and typography were repaired, preserving mathematical content and scope. This supplied editorial label does not establish mathematical correctness or completeness of the formulation.
\end{itemize}
\end{document}
